\chapter{Hardware}

La simulazione Wokwi descritta in \texttt{diagram.json} usa una Arduino Uno, un servo, pulsanti, LED, sensori e un LCD 16x2. Il cablaggio e coerente con i pin definiti nel firmware, cosi che la simulazione sia una rappresentazione fedele del comportamento atteso su hardware reale.

\section{Componenti}
\begin{itemize}
  \item Arduino Uno
  \item Servo per tornello
  \item 2 pulsanti (IN/OUT) con resistori di pull-down
  \item LED verde e rosso (semaforo)
  \item LED blu e arancione (raffreddamento/riscaldamento)
  \item LED bianco (umidita)
  \item Fotoresistenza (LDR)
  \item Sensore NTC per temperatura
  \item Sensore DHT22 per umidita
  \item 4 LED gialli per plafoniera (pilotati in parallelo)
  \item LCD 16x2
\end{itemize}

\section{Criteri di scelta}
I componenti sono stati selezionati per semplicita di integrazione e disponibilita di librerie stabili. L'NTC offre una misura analogica lineare per il termostato, mentre il DHT22 fornisce umidita con interfaccia digitale semplice. L'LCD 16x2 rappresenta un compromesso tra costo e leggibilita, con un'interfaccia standardizzata tramite \texttt{LiquidCrystal}.

\section{Mappa pin}
\begin{longtable}{@{}lll@{}}
\toprule
Funzione & Pin & Note \\
\midrule
Servo tornello & D3 & PWM \\
Pulsante IN & D6 & ingresso digitale \\
Pulsante OUT & D7 & ingresso digitale \\
LED verde semaforo & D4 & uscita digitale \\
LED rosso semaforo & D5 & uscita digitale \\
LED raffreddamento & D8 & uscita digitale \\
LED riscaldamento & D9 & uscita digitale \\
Plafoniera (PWM) & D10 & dimmer luci \\
LED umidita & D11 & uscita digitale \\
Sensore temperatura (NTC) & A0 & ingresso analogico \\
Fotoresistenza & A1 & ingresso analogico \\
LCD RS & D12 & \\
LCD EN & D13 & \\
LCD D4 & A2 & \\
LCD D5 & A3 & \\
LCD D6 & A4 & \\
LCD D7 & A5 & \\
DHT22 & D2 & pin definito in \texttt{humidity.h} \\
\bottomrule
\end{longtable}

\section{Schema grafico del cablaggio}
La figura seguente mostra una vista sintetica del cablaggio con Arduino al centro, sensori a sinistra e attuatori/display a destra. I collegamenti sono raggruppati per tipo di segnale (analogico, digitale, PWM, bus LCD) per migliorare la leggibilita mantenendo invariata la mappa pin firmware.

\begin{center}
\resizebox{\textwidth}{!}{%
\begin{tikzpicture}[>=latex, thick, node distance=1.2cm]
  \tikzstyle{arduino}=[draw, rounded corners, minimum width=4.2cm, minimum height=8.5cm, align=center, fill=gray!14]
  \tikzstyle{sensor}=[draw, rounded corners, minimum width=3.2cm, minimum height=0.9cm, align=center, fill=green!12]
  \tikzstyle{actor}=[draw, rounded corners, minimum width=3.5cm, minimum height=0.9cm, align=center, fill=blue!12]
  \tikzstyle{groupbox}=[draw, rounded corners, dashed, line width=0.8pt, fill=gray!3]

  \draw[red!70!black, line width=1.8pt] (-8.2,4.55) -- (8.2,4.55) node[right, black, font=\scriptsize\bfseries] {Rail +5V};
  \draw[black!70, line width=1.8pt] (-8.2,-4.55) -- (8.2,-4.55) node[right, black, font=\scriptsize\bfseries] {Rail GND};

  \node[arduino] (uno) at (0,0) {Arduino UNO\\\small Firmware Pinacoteca};

  \node[draw, fill=green!22, minimum width=0.42cm, minimum height=6.6cm, anchor=east] (hdrL) at ($(uno.west)+(0,0)$) {};
  \node[draw, fill=blue!22, minimum width=0.42cm, minimum height=6.6cm, anchor=west] (hdrR) at ($(uno.east)+(0,0)$) {};
  \node[rotate=90, font=\scriptsize\bfseries] at (hdrL.center) {HEADER INGRESSI};
  \node[rotate=90, font=\scriptsize\bfseries] at (hdrR.center) {HEADER USCITE};

  \node[groupbox, minimum width=3.8cm, minimum height=4.8cm, anchor=center] (sensorsbox) at (-6,0.6) {};
  \node[above left=0.0cm and -0.1cm of sensorsbox.north west, font=\footnotesize\bfseries] {Sensori / Ingressi};

  \node[groupbox, minimum width=4.1cm, minimum height=7.5cm, anchor=center] (actorsbox) at (6,0.1) {};
  \node[above left=0.0cm and -0.1cm of actorsbox.north west, font=\footnotesize\bfseries] {Attuatori / Display};

  \node[sensor] (ntc) at (-6,3.0) {NTC temperatura (A0)};
  \node[sensor] (ldr) at (-6,1.7) {Fotoresistenza LDR (A1)};
  \node[sensor] (dht) at (-6,0.4) {DHT22 umidita (D2)};
  \node[sensor] (btnin) at (-6,-1.2) {Pulsante IN (D6)};
  \node[sensor] (btnout) at (-6,-2.5) {Pulsante OUT (D7)};

  \node[actor] (servo) at (6,3.2) {Servo tornello (D3)};
  \node[actor] (tl) at (6,2.0) {Semaforo verde/rosso (D4,D5)};
  \node[actor] (clima) at (6,0.8) {LED clima cool/heat (D8,D9)};
  \node[actor] (light) at (6,-0.4) {Plafoniera PWM (D10)};
  \node[actor] (humled) at (6,-1.6) {LED umidita (D11)};
  \node[actor] (lcd) at (6,-3.1) {LCD 16x2 (D12,D13,A2,A3,A4,A5)};

  \node[font=\scriptsize\bfseries, text=green!50!black] at (-6,4.05) {LATO INGRESSI};
  \node[font=\scriptsize\bfseries, text=blue!50!black] at (6,4.05) {LATO USCITE};

  \draw[->, teal!80!black, line width=1.2pt] (ntc.east) -- node[above,sloped]{A0} ($(uno.west)+(0,2.8)$);
  \draw[->, teal!80!black, line width=1.2pt] (ldr.east) -- node[above,sloped]{A1} ($(uno.west)+(0,1.8)$);
  \draw[->, purple!75!black, line width=1.2pt] (dht.east) -- node[above,sloped]{D2} ($(uno.west)+(0,0.8)$);
  \draw[->, purple!75!black, line width=1.2pt] (btnin.east) -- node[below,sloped]{D6} ($(uno.west)+(0,-1.3)$);
  \draw[->, purple!75!black, line width=1.2pt] (btnout.east) -- node[below,sloped]{D7} ($(uno.west)+(0,-2.3)$);

  \draw[->, orange!90!black, line width=1.2pt] ($(uno.east)+(0,3.0)$) -- node[above,sloped]{D3 PWM} (servo.west);
  \draw[->, purple!75!black, line width=1.2pt] ($(uno.east)+(0,2.1)$) -- node[above,sloped]{D4,D5} (tl.west);
  \draw[->, purple!75!black, line width=1.2pt] ($(uno.east)+(0,1.1)$) -- node[above,sloped]{D8,D9} (clima.west);
  \draw[->, orange!90!black, line width=1.2pt] ($(uno.east)+(0,-0.2)$) -- node[above,sloped]{D10 PWM} (light.west);
  \draw[->, purple!75!black, line width=1.2pt] ($(uno.east)+(0,-1.3)$) -- node[above,sloped]{D11} (humled.west);
  \draw[->, blue!75!black, line width=1.2pt] ($(uno.east)+(0,-2.6)$) -- node[below,sloped]{Bus LCD D12,D13,A2-A5} (lcd.west);

  \fill[white, draw=black] ($(uno.west)+(0,2.8)$) circle (0.07);
  \fill[white, draw=black] ($(uno.west)+(0,1.8)$) circle (0.07);
  \fill[white, draw=black] ($(uno.west)+(0,0.8)$) circle (0.07);
  \fill[white, draw=black] ($(uno.west)+(0,-1.3)$) circle (0.07);
  \fill[white, draw=black] ($(uno.west)+(0,-2.3)$) circle (0.07);
  \fill[white, draw=black] ($(uno.east)+(0,3.0)$) circle (0.07);
  \fill[white, draw=black] ($(uno.east)+(0,2.1)$) circle (0.07);
  \fill[white, draw=black] ($(uno.east)+(0,1.1)$) circle (0.07);
  \fill[white, draw=black] ($(uno.east)+(0,-0.2)$) circle (0.07);
  \fill[white, draw=black] ($(uno.east)+(0,-1.3)$) circle (0.07);
  \fill[white, draw=black] ($(uno.east)+(0,-2.6)$) circle (0.07);

  \node[draw, rounded corners, fill=white, align=left, font=\footnotesize, anchor=north west] (legend) at (-8.3,-4.5) {
    \textbf{Legenda}\\
    \textcolor{teal!80!black}{\rule{0.8cm}{1.5pt}}\ Ingressi analogici\\
    \textcolor{purple!75!black}{\rule{0.8cm}{1.5pt}}\ I/O digitali\\
    \textcolor{orange!90!black}{\rule{0.8cm}{1.5pt}}\ Uscite PWM\\
    \textcolor{blue!75!black}{\rule{0.8cm}{1.5pt}}\ Bus parallelo LCD
  };
\end{tikzpicture}%
}
\end{center}

\section{Note di cablaggio}
\begin{itemize}
  \item I LED sono collegati con resistori da 220\,\textohm per limitare la corrente.
  \item I pulsanti usano resistori da 10\,k\textohm per pull-down, evitando letture flottanti.
  \item Il DHT22 usa il pin digitale 2 e richiede alimentazione a 5V.
  \item La plafoniera e pilotata in PWM sul pin D10, soluzione che semplifica la regolazione dell'intensita.
\end{itemize}
